\section{Introduction}
\label{sec:intro}

% --- 1 page max ---
% Story arc:
%   (a) HRP is widely used but its hard clustering and bisection are non-differentiable
%   (b) Differentiable optimization layers have unlocked end-to-end systems in other domains
%   (c) We propose DHRP: a soft-gating tree that generalizes HRP via differentiability
%   (d) E&D contribution: 8-universe benchmark with multi-seed PSR + factor alphas
%   (e) Honest negative finding: LLM augmentation fails under rigorous multi-seed evaluation

Portfolio optimization sits at the intersection of statistical estimation
and constrained decision-making. The classical mean-variance framework of
\citet{Markowitz1952} provides an elegant theoretical solution but is
notoriously sensitive to estimated inputs: small errors in expected
returns or covariances are amplified by the optimization, leading to
unstable allocations that perform poorly out of sample
\citep{DeMiguel2009,Ledoit2004}. Hierarchical Risk Parity (HRP),
introduced by \citet{LopezdePrado2016}, sidesteps this fragility by
working with covariance information alone and routing capital down a
binary tree built from agglomerative clustering of asset correlation
distances. HRP produces stable, diversified portfolios and has become a
practitioner standard, with subsequent work exploring linkage variants
\citep{Raffinot2018} and risk-budgeting extensions \citep{Roncalli2013}.

From a machine-learning perspective, HRP has a critical limitation: its
clustering and recursive bisection steps are non-differentiable, so it
cannot be jointly trained with upstream feature extractors. Recent
advances in differentiable optimization layers
\citep{Amos2017,Agrawal2019,Bolte2021,Pan2024BPQP} have shown that
embedding structured decision-making within neural networks unlocks
end-to-end systems that outperform two-stage ``predict then optimize''
pipelines \citep{Donti2017,Wilder2019,Elmachtoub2022SPO}. Concurrently,
soft decision trees \citep{FrosstHinton2017,Popov2020NODE,Mao2025DiffTree}
and gated mixture-of-experts \citep{Shazeer2017} have demonstrated that
hierarchical routing can be made differentiable while retaining the
inductive bias that makes trees interpretable.

This paper makes both architectures available to portfolio allocation.
We replace HRP's clustering and bisection with a temperature-scaled
soft-gating tree that recovers classical HRP as routing temperature
$\tau \to 0$, and we evaluate the resulting layer (DHRP) against ten
strong baselines across eight asset universes with rigorous multi-seed
reporting (10 seeds, mean$\pm$std + Probabilistic Sharpe Ratios).
We further test whether augmenting DHRP with FinBERT \citep{Huang2023}
embeddings of asset-attributed financial news helps---and find,
honestly, that it does not, joining a growing benchmark literature
\citep{Tan2024AreLMs,Chen2025StockBench,Xie2024FinBen} cautioning
against single-seed claims of LLM benefit in finance.

\paragraph{Contributions.} This paper makes three contributions:
\begin{enumerate}
  \item \textbf{DHRP architecture (Section~\ref{sec:method}).} A
    temperature-scaled soft-gating tree that recovers classical HRP as
    routing temperature $\tau \to 0$ (Proposition~\ref{prop:recovery}),
    enabling end-to-end gradient flow through hierarchical allocation.
  \item \textbf{DHRP-8 benchmark (Section~\ref{sec:benchmark}).} An
    8-universe portfolio evaluation framework (developed markets, emerging
    markets, commodities, sectors, global multi-asset, factors, crypto,
    bonds) with eleven methods, multi-seed ($n=10$) reporting,
    Probabilistic Sharpe Ratios, Fama--French and AQR commodity factor
    alphas, transaction-cost sensitivity, and Croissant~1.1 metadata for
    full reproducibility.
  \item \textbf{Empirical findings (Sections~\ref{sec:results},
    \ref{sec:llmablation}).} DHRP achieves statistically significant
    alpha in commodities and crypto with PSR $> 0.99$. Mean-variance
    dominates equity-heavy universes; no method universally dominates
    (SPA $p=0.30$). LLM augmentation provides no statistical benefit
    under multi-seed reporting and regresses in commodities, contrary to
    single-seed claims in the literature.
\end{enumerate}

% Roadmap paragraph
The remainder of the paper is organized as follows.
Section~\ref{sec:related} surveys related work in differentiable
portfolio optimization, hierarchical methods, and LLM-finance.
Section~\ref{sec:method} formalizes the DHRP architecture and its loss.
Section~\ref{sec:benchmark} describes the 8-universe benchmark protocol.
Section~\ref{sec:results} reports headline results.
Section~\ref{sec:llmablation} documents the LLM negative ablation.
Section~\ref{sec:limits} discusses limitations.
